\documentclass[11pt,a4paper]{article}
\usepackage[T1]{fontenc}
\usepackage{lmodern,amsmath,amssymb,amsthm,mathtools,booktabs,array,longtable,microtype}
\usepackage[margin=27mm]{geometry}
\usepackage[colorlinks=true,allcolors=blue]{hyperref}
\newtheorem{theorem}{Theorem}[section]
\newtheorem{proposition}[theorem]{Proposition}
\newtheorem{lemma}[theorem]{Lemma}
\newtheorem{corollary}[theorem]{Corollary}
\theoremstyle{remark}\newtheorem{remark}[theorem]{Remark}
\newcommand{\U}{\mathrm U}
\newcommand{\SU}{\mathrm {SU}}
\newcommand{\PU}{\mathrm {PU}}
\newcommand{\su}{\mathfrak {su}}
\newcommand{\slie}{\mathfrak {sl}}
\newcommand{\C}{\mathbb C}
\newcommand{\Tr}{\operatorname{Tr}}
\newcommand{\diag}{\operatorname{diag}}
\newcommand{\Span}{\operatorname{span}}
\newcommand{\Cl}{\operatorname{Cliff}}

\newcommand{\A}{\mathcal A}
\newcommand{\B}{\mathcal B}
\newcommand{\h}{\mathfrak h}
\newcommand{\one}{\mathbf 1}
\newcommand{\mc}{\operatorname{mc}}
\newcommand{\MC}{\operatorname{MC}}
\newcommand{\Conf}{\operatorname{Conf}}
\newcommand{\Rep}{\operatorname{Rep}_{\mathrm{fd}}}
\newcommand{\im}{\operatorname{im}}
\newcommand{\Id}{\operatorname{Id}}
\title{Generalized Burau representations and quantum gates:\\arrangements, compact closures, and prior results}
\author{Haru Negami}
\date{Note, version 2, 10 November 2026}
\hypersetup{pdfauthor={Haru Negami},pdftitle={Generalized Burau representations and quantum gates}}
\begin{document}
\maketitle
\begin{abstract}
We study two constructions of Burau-type representations and their
quantum-gate realizations. For finite Coxeter arrangements, we collect
Squier's explicit matrices and determine when their invariant Hermitian
forms are positive definite in terms of the Coxeter exponents. A direct
rank-two proof gives a compact-closure formula for
$0<|\theta|<\min\{\pi/h,\pi/10\}$, including roots of unity.
For $E_6,E_7,E_8$, this formula holds throughout the parameter interval
where the invariant Hermitian form is positive definite, except at the
undeformed Coxeter specialization. These exceptional-type conclusions
also follow from Freedman--Larsen--Wang and are not claimed as new
classification theorems. The eight-dimensional gates give the stated
non-Clifford and three-qubit universality conclusions.

For arbitrary affine arrangements, we now carry out the separate
construction from a unitary rank-one seed through directional middle
convolution to quantum gates. Specializing Negami's prior invariant-form
and signature theorems gives explicit definiteness conditions, including
the middle quotient at resonance. We verify colored braid transport for
unequal weights, give a sufficient condition for full unitary closure
and non-Clifford gates, and an exact finite-image test for rational
phases with finite-image base character. The output lives on the
directional closure when collision hyperplanes must be added.
The appendices retain the Dynkin dictionary, arbitrary-rank algebraic
convolution formulas, and the based six-line example. This note provides
explicit applications, proofs, conventions, and reproducible code; the
two constructions are not identified in general.
\end{abstract}

\section{Introduction and statement}\label{sec:general-burau}
In earlier work, we proposed a method for constructing quantum gates
from monodromy representations obtained by applying middle convolution
to KZ-type systems \cite{NegamiGates,NegamiKZ}.
The arrangement construction considered here follows that method:
start with a unitary rank-one local system on a general affine hyperplane
arrangement, apply directional middle convolution, and use the author's
invariant Hermitian form and signature results
\cite{NegamiKZ,NegamiSignature} to determine when the resulting matrices
can be realized as unitary gates.

Two families are treated with their respective domains and hypotheses.
Sections~1--6 retain Squier's representations of Artin--Tits groups
\cite{Squier}, their positive-definite parameter ranges in finite Coxeter
type, and the corresponding closure and gate results.
Section~\ref{sec:mc-gates}, added in version~2, carries out the
arrangement-to-MC-to-gates construction itself, including the actual
middle quotient and the base-group action. This construction uses
Adachi--Hiroe's directional closure \cite{AH,Hiroe} and the
Katz--Long--Moody framework \cite{HN}. The Burau connection is the
equal-weight homological specialization; it does not identify the
general MC output with every Coxeter-type representation.

Our emphasis is on explicit parameter conditions, on the invariant form
preserved by the complete monodromy, and on the distinction between
finite-order generators and a finite generated group. The multiplicative
MC construction in Section~\ref{sec:mc-gates} also distinguishes
resonant middle quotients from unrestricted additive specializations.
The arbitrary-rank additive formulas and the earlier fixed-meridian
calculations remain in the appendices.

\paragraph{General Coxeter type versus arbitrary arrangements.}
For any finite-rank Coxeter matrix, not necessarily of finite type, the
known construction on $\C(t)^r$ is
\begin{equation}\label{eq:coxeter-burau}
 R_i e_j=\begin{cases}
 -t^2 e_i,&j=i,\\
 e_j+2t\cos(\pi/m_{ij})e_i,&j\ne i,
 \end{cases}\qquad \cos(\pi/\infty)=1.
\end{equation}
The matrices satisfy the Artin relations, have determinant $-t^2$, and
satisfy $(R_i-I)(R_i+t^2I)=0$; they specialize for every $t\ne0$
\cite{Squier,BQ}.  Our parameter is the negative of that in
\cite[Definition 2.1]{BQ}, and agrees with \cite{SageArtin}.
An arbitrary affine arrangement need not have a Coxeter matrix or a
Dynkin type.  Appendix~\ref{sec:dynkin} explains the type notation;
Section~\ref{sec:mc-gates} and
Appendices~\ref{sec:general-arrangement}--\ref{sec:descent} treat a
different, direction-dependent construction for arbitrary arrangements.
No identification of the exceptional Coxeter representations with a
middle-convolution output is assumed.

The closure results below require finite irreducible Coxeter type.

Let $W$ be a finite irreducible Coxeter group of rank $r\geq2$, with
Coxeter matrix $(m_{ij})$, Coxeter number $h$, and Artin group $B_W$.
Write $P_W=\ker(B_W\to W)$.  Define the real symmetric matrix
\begin{equation}\label{eq:A}
 A_{ii}=0,\qquad A_{ij}=2\cos(\pi/m_{ij})\quad(i\ne j).
\end{equation}
For $t=e^{i\theta}$ put
\begin{equation}\label{eq:H}
 H_\theta=2\cos\theta\,I-A,
 \qquad R_i=I-t e_i e_i^*H_\theta.
\end{equation}
These are the generalized Burau matrices, with
$R_ie_i=-t^2e_i$ and $R_ie_j=e_j+tA_{ij}e_i$ for $j\ne i$.
When $H_\theta>0$, let
\begin{equation}\label{eq:U}
 S=H_\theta^{1/2},\qquad U_i=SR_iS^{-1},\qquad
 \lambda=-e^{2i\theta},\qquad D_\theta=\overline{\langle\lambda\rangle}
 \subset\U(1).
\end{equation}
Closures throughout the paper are in the ordinary compact-group topology.

\begin{theorem}[Explicit closure formula]\label{thm:main}
Suppose
\begin{equation}\label{eq:range}
 0<|\theta|<\min\{\pi/h,\pi/10\}.
\end{equation}
Then the unitarized representation satisfies
\begin{align}
 \overline{\rho_\theta(B_W)}
 &=\{g\in\U(r):\det g\in D_\theta\},\label{eq:closure}\\
 \overline{\rho_\theta(P_W)}
 &=\{g\in\U(r):\det g\in
            \overline{\langle\lambda^2\rangle}\}.\label{eq:pureclosure}
\end{align}
In particular, both projective images are dense in $\PU(r)$.
If $\theta/\pi$ is irrational, both closures are $\U(r)$.
If $\lambda$ has finite order $k$, the closures have respectively
$k$ and $k/\gcd(k,2)$ connected components, each a coset of $\SU(r)$.

For $h\geq10$, the hypothesis covers the full positive interval
$0<|\theta|<\pi/h$.  In $|\theta|<\pi/h$ the Artin image is then finite
if and only if $\theta=0$; at $0$ the image is the reflection representation
of $W$, and the pure image is trivial.
\end{theorem}

The exceptional types have $h=12,18,30$, respectively.  Thus the theorem
is exhaustive within their positive intervals.  It is not a classification
of specializations outside those intervals, or at degenerate endpoints.

\paragraph{Relation to prior work.}
The representation and form are due to Squier \cite{Squier}.
The exceptional-type conclusion of Theorem~\ref{thm:main} is a
consequence of Freedman--Larsen--Wang \cite{FLW}; the precise reduction
is given in Proposition~\ref{prop:flw}.  Thus an explicit interval or a
root-of-unity specialization does not, by itself, establish novelty.
Marin's generic compact specialization theorem \cite[Theorem~4]{Marin}
and the rank-two analysis of Funar--Kohno \cite{FK} provide further
precedent.  The finite $W(E_8)$ connection with three-qubit gates is also
known \cite{Planat,ARW}.  Selective-phase gate descriptions and the Toffoli
comparison overlap with the author's separate working draft
\cite[Section~6]{Negami}.  Section~\ref{sec:overlap} distinguishes known
constructions, consequences of existing theorems, and explicit calculations.

\section{Invariant forms and explicit unitary gates}
The matrix in \eqref{eq:H} is the specialization of Squier's two-variable
form obtained by setting both variables equal to $t$ and dividing the
Gram matrix by $t$.  The matrices $R_i$ satisfy the Artin relations
\cite{Squier}.  In our convention, type $A_r$ corresponds to reduced
Burau with classical parameter $q=t^2$.

\begin{proposition}\label{prop:form}
For $|t|=1$, $R_i^*H_\theta R_i=H_\theta$.  If $d_1,\ldots,d_r$
are the Coxeter exponents, the eigenvalues of $H_\theta$ are
\begin{equation}\label{eq:spectrum}
 2\cos\theta-2\cos(\pi d_j/h),\qquad 1\leq j\leq r.
\end{equation}
For $-\pi<\theta\leq\pi$, this particular form is positive definite
exactly when $|\theta|<\pi/h$.
\end{proposition}
\begin{proof}
Set $E_i=e_ie_i^*$.  Expansion gives
\[
 R_i^*H_\theta R_i-H_\theta
 =[-t-\bar t+|t|^2(H_\theta)_{ii}]H_\theta E_iH_\theta=0.
\]
The spectrum of $A$ is $2\cos(\pi d_j/h)$; for crystallographic types
this is the standard Cartan/adjacency spectral identity, with diagonal
symmetrization in the nonsimply-laced cases \cite{Damianou}.
The finite Coxeter version uses the same geometric Gram matrix.
The smallest exponent is $1$, so the smallest eigenvalue of $H_\theta$
is $2\cos\theta-2\cos(\pi/h)$.
\end{proof}

At a degenerate parameter, $\ker H_\theta$ is fixed pointwise by all
$R_i$ and the form descends to a nondegenerate form on the quotient.
We do not use that quotient in Theorem~\ref{thm:main}.  The negative-definite
region can also be unitarized after changing the overall sign; our
finite-image statement is explicitly restricted to the positive interval.

In the positive interval set
\begin{equation}\label{eq:rankone}
 p_i=\frac{Se_i}{\sqrt{2\cos\theta}},\qquad
 U_i=I+(\lambda-1)p_ip_i^*,\qquad
 \langle p_i,p_j\rangle=-\frac{A_{ij}}{2\cos\theta}\quad(i\ne j).
\end{equation}
The $p_i$ form a basis of unit vectors.  Each $U_i$ is the identity on
$p_i^\perp$ and multiplies $p_i$ by $\lambda$; hence
\begin{equation}\label{eq:trace}
 \det U_i=\lambda,\qquad \Tr U_i=r-1+\lambda.
\end{equation}
The same positive square root $S$ is used for every generator.

For $E_r$, $r=6,7,8$, use the edges
\[
 \mathcal E_r=\{\{i,i+1\}:1\leq i\leq r-2\}\cup\{\{3,r\}\}.
\]
Then $A$ is the ordinary adjacency matrix.  This convention fixes all
matrices without a choice of root coordinates.
\begin{center}
\begin{tabular}{ccll}
\toprule
Type & $h$ & Exponents & Positive interval\\
\midrule
$E_6$ & $12$ & $1,4,5,7,8,11$ & $|\theta|<\pi/12$\\
$E_7$ & $18$ & $1,5,7,9,11,13,17$ & $|\theta|<\pi/18$\\
$E_8$ & $30$ & $1,7,11,13,17,19,23,29$ & $|\theta|<\pi/30$\\
\bottomrule
\end{tabular}
\end{center}
For example, $U_3=I+(\lambda-1)(Se_3)(Se_3)^*/(2\cos\theta)$.
In the original basis, row $3$ of $R_3$ has entries $t$ in columns
$2,4,r$, entry $-t^2$ in column $3$, and zeros elsewhere; all other
rows are those of the identity.

\section{From two-dimensional subspaces to the full closure}\label{sec:direct}
The proof below does not infer density from irreducibility or from
infinite order alone.

\begin{lemma}\label{lem:two}
Let $p,q$ be independent unit vectors with
$0<|\langle p,q\rangle|<1$.  On $T=\Span\{p,q\}$ consider
$V_p=I+(\lambda-1)pp^*$ and $V_q=I+(\lambda-1)qq^*$, where
$\lambda=-e^{2i\theta}$ and $0<|\theta|<\pi/10$.
Their projective image is dense in $\PU(T)$, and their unitary closure
contains $\SU(T)$.
\end{lemma}
\begin{proof}
Under $\PU(2)\simeq\mathrm{SO}(3)$, each generator is a rotation of
angle $\alpha=\pi-2|\theta|\in(4\pi/5,\pi)$.  Their axes are distinct:
equality of the unoriented axes would mean that the two eigenlines are
equal or orthogonal, both excluded by hypothesis.

The finite rotation groups are cyclic, dihedral, tetrahedral, octahedral,
and icosahedral.  In the first two, every rotation with angle different
from $\pi$ has the same axis.  In the last three, every angle smaller
than $\pi$ is at most $4\pi/5$.  Thus the projective group is infinite.
An infinite proper closed subgroup of $\mathrm{SO}(3)$ has identity
component a circle and is contained in its normalizer.  In that normalizer
all rotations other than half turns have the circle's axis, again a
contradiction.  The projective closure is therefore $\mathrm{SO}(3)$.

If $\mathfrak k\subset\mathfrak u(2)$ is the Lie algebra of the unitary
closure, its projection onto $\mathfrak{pu}(2)$ is surjective.  Hence
$[\mathfrak k,\mathfrak k]\subset\su(2)$ projects onto $\su(2)$ and
equals it.  The closure contains $\SU(2)$.
\end{proof}

\begin{lemma}\label{lem:fusion}
Let $v_1,\ldots,v_r$ be a basis of a Hilbert space $V$, and let $\Gamma$
be a connected graph on these indices.  A closed subgroup of $\U(V)$
containing $\SU(\Span\{v_i,v_j\})$, acting trivially on its orthogonal
complement, for every edge $\{i,j\}$ contains $\SU(V)$.
\end{lemma}
\begin{proof}
We use a connected ordering of a spanning tree.  Suppose $\SU(S)$ has
already been obtained and an edge supplies $\SU(T)$, where $\dim T=2$
and $S\cap T=\C u$ for a unit vector $u$.  Choose a unit vector
$w\in T\cap u^\perp$ and write $w=a+bz$, with $a\in S\cap u^\perp$,
$z\perp S$ a unit vector, and $b\ne0$.
In the complexified generated Lie algebra, $uw^*$ belongs to
$\slie(T)$ and $ua^*$ to $\slie(S)$.  Subtraction gives $uz^*$;
similarly one obtains $zu^*$.  Commutators with $\slie(S)$ give
$xz^*,zx^*$ for all $x\in S$, and their commutators give the missing
trace-zero diagonal operator.  They generate $\slie(S+T)$.
The original compact Lie algebra therefore contains $\su(S+T)$.
Induction, starting from one edge, gives the assertion.
\end{proof}

\begin{proof}[Proof of Theorem~\ref{thm:main}]
For an edge of the Coxeter graph, \eqref{eq:rankone} and positivity give
$0<|\langle p_i,p_j\rangle|<1$.  The two gates preserve their span and
fix its orthogonal complement pointwise.  Lemma~\ref{lem:two} gives an
embedded $\SU(2)$ in the closure on each edge.  Connectedness and
Lemma~\ref{lem:fusion} imply that the full closure $G$ contains $\SU(r)$.

The determinant of every word is a power of $\lambda$.  By compactness,
$\det G=\overline{\langle\lambda\rangle}=D_\theta$.  Since the kernel
of determinant is contained in $G$, this proves \eqref{eq:closure}.

The pure group has finite index in the Artin group.  Its image closure
$G_P$ consequently has finite index in $G$, so contains $G^\circ$ and
in particular $\SU(r)$.  The exponent-sum homomorphism of the Artin group
is even on $P_W$, since its parity is the Coxeter sign character.  Moreover
$\sigma_i^2\in P_W$ has exponent sum $2$.  Hence the determinant image
of $P_W$ is exactly $\langle\lambda^2\rangle$, proving
\eqref{eq:pureclosure}.  The assertions about rational and irrational
parameters follow from closed subgroups of $\U(1)$.

At $\theta=0$, the matrices $R_i$ are the faithful geometric reflection
representation of $W$.  Every other parameter in \eqref{eq:range}
has closure containing the positive-dimensional group $\SU(r)$ and
therefore has infinite image.  If $h\geq10$, the positive interval is
contained in $|\theta|<\pi/10$.
\end{proof}

\begin{remark}[Scope for other types]
For crystallographic groups, the full-positive-interval conclusion also
applies to $A_n$ with $n\geq9$, $B_n=C_n$ with $n\geq5$, $D_n$ with
$n\geq6$, and $F_4$.  Smaller Coxeter numbers are still covered on the
smaller interval in \eqref{eq:range}; no assertion is made here about
their remaining positive parameters.  Merely assuming that $\lambda$
is a root of unity would not prove finite image.
\end{remark}

\section{Clifford obstruction and quantum universality}
We fix the usual identification $\C^8=(\C^2)^{\otimes3}$ in the
orthonormal coordinates of \eqref{eq:U}.  The Clifford group means the
normalizer of the tensor-product Pauli group, with arbitrary global
phases allowed.  Its binary symplectic action is standard \cite{DDM}.

\begin{lemma}\label{lem:clifford}
If $C$ is a nonscalar Clifford unitary on $n$ qubits and $d=2^n$, then
\begin{equation}\label{eq:tracebound}
 |\Tr C|^2\leq d^2/2.
\end{equation}
Consequently, for $n\geq3$ and $\lambda\ne1$, no unitary with spectrum
$\{1^{(d-1)},\lambda\}$ is Clifford, in any orthonormal encoding.
\end{lemma}
\begin{proof}
Let $F\in\mathrm{Sp}(2n,2)$ be the induced permutation of Pauli labels.
The conjugation operator $\operatorname{Ad}_C$ is a signed permutation
on the Hermitian Pauli basis, and
$\Tr(\operatorname{Ad}_C)=|\Tr C|^2$.
If $F\ne I$, at most $2^{2n-1}$ labels are fixed, so its trace is at most
that number.  If $F=I$, the Clifford is a Pauli up to phase; unless
scalar, its trace is zero.  This proves \eqref{eq:tracebound}.
For the stated spectrum,
$|\Tr U|\geq d-2$, whereas $(d-2)^2>d^2/2$ for $d\geq8$.
Trace magnitude is unchanged by a global phase or a common unitary
change of encoding.
\end{proof}

\begin{corollary}\label{cor:e8}
For type $E_8$ and $|\theta|<\pi/30$, every generator $U_i$ is
non-Clifford.  For $0<|\theta|<\pi/30$, the gates
$\{U_i,U_i^{-1}:1\leq i\leq8\}$ are approximately universal on three
qubits, up to global phase.  At $\theta=0$ all generators remain
non-Clifford, but the generated group is the finite group $W(E_8)$.
\end{corollary}
\begin{proof}
Here $\lambda\ne1$ and $\Tr U_i=7+\lambda$, so
$|\Tr U_i|^2\geq36>32$.  Apply Lemma~\ref{lem:clifford} and the
projective closure in Theorem~\ref{thm:main}.
\end{proof}

Non-Cliffordness alone is thus insufficient for universality, even within
this single continuous family.  A fixed three-qubit gate set is meant
here.  If its gates can be applied to arbitrary selected triples in a
larger register, one obtains circuit universality by approximating arbitrary
one- and two-qubit gates inside a triple \cite{Barenco}.  No implementation cost or
fault-tolerance advantage is asserted.

\paragraph{An exact root-of-unity example.}
Take $\theta=\pi/60$.  Then
\[
 t=e^{i\pi/60},\qquad \lambda=e^{2\pi i\,31/60},\qquad U_i^{60}=I.
\]
Although all eight gates have exact order $60$, their group is infinite,
its projective closure is $\PU(8)$, and its full closure is
\[
 \{g\in\U(8):\det g\in\mu_{60}\}.
\]
The pure-group closure has determinant in $\mu_{30}$.  Inverses are
already positive powers of the gates in this example.

\paragraph{Gate realization.}
If a unitary $V_i$ satisfies $V_i|111\rangle=p_i$, then
\begin{equation}\label{eq:circuit}
 U_i=V_i\diag(1,1,1,1,1,1,1,\lambda)V_i^*.
\end{equation}
Thus each generator is a conjugated three-qubit controlled phase.
Such a $V_i$ can be constructed by completing $p_i$ to an orthonormal
basis, or by a Householder map.  Formula \eqref{eq:circuit} is a matrix
realization, not a claim that $V_i$ is a free or Clifford operation.
At $\theta=0$ it is a conjugate of CCZ, hence also of Toffoli.
This is an individual conjugacy statement, not a simultaneous
identification of the entire generator set with one gate.

\paragraph{The other exceptional dimensions.}
Types $E_6,E_7$ give projectively universal gates on $\C^6,\C^7$.
One may use $\C^6\simeq\C^2\otimes\C^3$, or encode either space in
three qubits.  Under the latter choice, the extension
$U_i\oplus I_{8-r}$ is still a nontrivial rank-one phase gate and is
non-Clifford by Lemma~\ref{lem:clifford}.  The group generated by these extensions preserves
the code and the unused subspace, so is not universal on all of $\C^8$.

\section{Arrangement interpretation and reproducibility}
Let $X_W$ be the complement of the complexified reflection arrangement.
Then $P_W=\pi_1(X_W)$ and $B_W=\pi_1(X_W/W)$.  Formula
\eqref{eq:pureclosure} therefore gives the compact closure for the
actual arrangement-complement representation, not only for the quotient
by $W$.  For a general arrangement and a map $f$ into $X_W$,
a pullback representation inherits the form, but its image can
be a much smaller subgroup.  Density does not follow without information
about $f_*(\pi_1)$.

An arbitrary arrangement need not have a Dynkin type.  The results above
concern reflection arrangements and do not supply a canonical universal
gate construction for all arrangements.  We also do not identify these
representations with a Katz--Long--Moody quotient; such a comparison is
unnecessary for the Coxeter closure proof. The direct MC gate
construction is instead developed in Section~\ref{sec:mc-gates}.

The accompanying \texttt{verify\_coxeter\_gates.py} constructs the
matrices and exports a concrete $E_8$ gate set at $\theta=\pi/60$.
It verifies the Artin and invariant-form identities exactly over
$\mathbb Z[t,t^{-1}]$ for $E_6,E_7,E_8$, and checks the positive
square-root construction, spectral formulas, determinants, and Clifford
trace obstruction numerically.  The numerical checks are consistency
checks; the density, finite-image classification, and non-Clifford
conclusions follow from the proofs above rather than finite word searches.


\section{Comparison with existing theorems}\label{sec:overlap}
The following reduction is included to make the overlap a mathematical
statement, rather than an inference from similar titles or abstracts.

\begin{proposition}[Exceptional-type closure as an existing-theorem consequence]
\label{prop:flw}
For $W=E_r$, $r=6,7,8$, and $0<|\theta|<\pi/h$, the inclusion
$\SU(r)\subset\overline{\rho_\theta(B_W)}$ follows from
\cite[Theorems~1.1, 1.4 and 1.6]{FLW}, without using the rank-two
density proof of Section~\ref{sec:direct}.
\end{proposition}
\begin{proof}
Put $G=\overline{\langle U_1,\ldots,U_r\rangle}$.  This is a compact
matrix group acting faithfully.  It acts irreducibly: a nonzero invariant
subspace contains some $p_i$, since the $p_i$ span and
$(U_i-I)v=(\lambda-1)p_i\langle p_i,v\rangle$.
Nonzero overlaps along the connected Coxeter diagram then put every
$p_j$ in that subspace.

All Artin generators in a connected simply-laced diagram are conjugate:
if $aba=bab$, then $(ab)a(ab)^{-1}=b$.
Consequently the conjugacy class of $U_1$ generates $G$ topologically.
Its eigenvalues are $1$ and $\lambda$ with multiplicities $r-1$ and $1$,
and $\lambda\ne\pm1$.  These are exactly the two-eigenvalue hypotheses
of \cite{FLW}.

If $G/Z(G)$ were finite, \cite[Theorem~1.6]{FLW} would imply
$U_1^m\in Z(G)$ for some $m\in\{3,4,5\}$.
By irreducibility this forces $\lambda^m=1$.  But the smaller angle
of $\lambda$ is $\pi-2|\theta|\in(4\pi/5,\pi)$, which excludes those
orders.  Thus $G/Z(G)$ has positive dimension.

Let $K=[G^\circ,G^\circ]$.  Applying
\cite[Theorems~1.1 and 1.4]{FLW}, $K$ acts irreducibly,
$G=G^\circ Z(G)$, and its representation belongs to their list.
In dimensions $6,7,8$, the alternatives to the defining or dual
representation of $\SU(r)$ all preserve a nondegenerate complex
bilinear form.  Explicitly, dimension $6$ also permits
$\bigwedge^2\C^4$ for $\SU(4)$ and the defining $\mathrm{Sp}(6)$
module; dimension $7$ has no alternative; dimension $8$ also permits
the defining $\mathrm{Sp}(8)$ module, the spin module for
$\mathrm{Spin}(7)$, and the vector and half-spin modules for
$\mathrm{Spin}(8)$.  Here $\mathrm{Sp}(d)$ denotes the compact
symplectic group in complex dimension $d$.

If such a bilinear form $B$ existed, the center of $G^\circ$ would
act by scalars, and $G=G^\circ Z(G)$ would imply
$U_i^{\mathsf T}BU_i=cB$ for some $c\ne0$.
The eigenvalue multiset would then be invariant under $z\mapsto c/z$.
The multiplicities $r-1$ and $1$ differ, so $1$ must map to itself and
$\lambda$ to itself.  Thus $c=1$ and $\lambda^2=1$, a contradiction.
Only the defining $\SU(r)$ image remains, as required.
\end{proof}

The determinant and finite-index arguments in the proof of
Theorem~\ref{thm:main} now give the full and pure closure formulas.
The exceptional-type finite-image statement and projective universality
are therefore further consequences, not independent classification
results.  This reduction does not assert that the same formulas occur
verbatim in \cite{FLW}.  For nonsimply-laced types, one cannot use the
single-conjugacy-class step without further justification; the direct
proof in Section~\ref{sec:direct} still applies in its stated interval.

\subsection*{Claim-by-claim status}
\renewcommand{\arraystretch}{1.13}
\begin{longtable}{@{}p{0.26\textwidth}p{0.34\textwidth}p{0.34\textwidth}@{}}
\toprule
Item in this note & Prior result or deduction & Status used here\\
\midrule\endhead
Generalized Burau matrices and invariant form & Squier \cite{Squier}; modern convention in \cite{BQ}. & Known construction, recalled explicitly.\\
Dynkin equations, ranks, exponents, and positivity & Root-system data and the Cartan spectral identity \cite{Etingof,Damianou}. & Standard data and a direct spectral consequence.\\
$E_6,E_7,E_8$ compact closure & Proposition~\ref{prop:flw}, applying \cite{FLW}. & Existing-theorem consequence; the rank-two proof is an alternative derivation.\\
Full unitary components and pure-group closure & Determinant image and finite index, once $\SU(r)$ is contained. & Elementary consequences of the closure statement.\\
Finite images in the positive $E$ interval & Infinite projective closure for $\theta\ne0$; Coxeter image at $0$. & A consequence within the stated interval, not a general classification.\\
Non-Clifford generators & Pauli/symplectic formalism \cite{DDM} and the trace estimate in Lemma~\ref{lem:clifford}. & Elementary consequence. An identical published formulation has not been located; no priority claim is made.\\
Three-qubit universality and finite $E_8$ limit & Closure implies universality; finite $W(E_8)$ gate realizations precede this note \cite{Planat,ARW}. & Application and explicit family; not a new universality criterion.\\
General-arrangement convolution and middle quotient & Existing constructions \cite{AH,Hiroe,HN}. & Recalled formulas with hypotheses; not a new functor.\\
Rank-one MC form and exact signature & Specializations of the author's results \cite{NegamiKZ,NegamiSignature}, with colored transport verified in Section~\ref{sec:mc-gates}. & Application to the complete arrangement monodromy, including its middle quotient.\\
MC unitary closure and finite images & Projector Lie-algebra generation and the cyclotomic argument of \cite{NegamiFinite}. & Sufficient irrational-phase density criterion; rational finite-image criterion under the stated base-character hypothesis.\\
Six-line monodromy and descent calculation & Burau pullback, fixed based loops, and a routine module quotient. & Explicit computation. Independent novelty of the particular presentation is unestablished.\\
\bottomrule
\end{longtable}

\paragraph{Parameter and scope checks.}
Marin's compact specialization result assumes a transcendental parameter
on the unit circle sufficiently near $1$ \cite[Theorem~4]{Marin}; it
cannot alone be quoted for every algebraic parameter in our interval.
That limitation does not rescue novelty of the exceptional-type result,
because Proposition~\ref{prop:flw} supplies a different existing theorem.
Funar--Kohno analyze rank-two root-of-unity images \cite{FK}; parameter
sign conventions must be matched before transferring their exceptional
orders.  A general arrangement pullback need not have the same image
as the whole braid group, as already noted above.

\paragraph{Research status.}
This is an expository and computational note.
The explicit gate matrices, the common unitarization, and the direct
closure proof are useful to retain, but do not by themselves demonstrate
a new principal theorem.  The literature comparison is targeted to the
claims made here, rather than a claim of exhaustive priority clearance.
No new theorem is inferred merely from failure to locate an identical
wording, parameterization, or example.

\section{Quantum gates from rank-one middle convolution}
\label{sec:mc-gates}
This section carries out the construction
\[
 (\A,\chi)\longmapsto\MC^Y_\lambda(\chi)
 \longmapsto H_{\MC}\longmapsto\{U_\gamma\}.
\]
Here $\lambda$ denotes the convolution parameter and is independent of
the Coxeter parameter $\theta$ and the notation $\lambda=-e^{2i\theta}$
used in Sections~1--6. The eigenangles below refer to multiplicative
monodromy; they are not the additive residues in the appendices.
In particular, the clockwise fiber convention here and the
counterclockwise pure-braid convention of Appendix~\ref{sec:six-lines}
must be matched before comparing matrix entries.

\subsection{Construction and its domain}
Let $\A$ be a finite affine hyperplane arrangement in $\C^d$ and choose
a one-dimensional direction $Y$. Write coordinates as $(x,y)$ and the
hyperplanes transverse to $Y$ as $y=f_i(x)$, $1\le i\le n$, where the
$f_i$ are distinct affine functions. Hyperplanes parallel to $Y$ have
equations depending only on $x$.
Let $\B=\A^Y$ be obtained by adjoining the inverse images of all base
hyperplanes $f_i=f_j$ for nonconstant differences. This is the directional
closure used in the arrangement middle convolution of Adachi--Hiroe
\cite{AH} and its Lie-algebraic realization \cite{Hiroe}.
The projection
\[
 M(\B)\longrightarrow X
\]
to the complement of the parallel and collision hyperplanes in the base
is a bundle of $n$-punctured affine lines. It has a continuous section:
one can choose a positive real $y$ larger than $1+\sum_i|f_i(x)|$.
Consequently its fundamental group contains the fiber group $F_n$ and
has a splitting $F_n\rtimes\pi_1(X)$. Fix such a splitting and compatible
based paths. The labelled root map $x\mapsto(f_1(x),\ldots,f_n(x))$
induces a pure braid action on the fiber.

Choose a unitary rank-one seed character $\chi$ on $\pi_1(M(\A))$,
pull it back to $M(\B)$, and write
\[
 q_i=\chi(x_i)=e^{2\pi i a_i},\qquad 0\le a_i<1,
 \qquad \lambda=e^{2\pi i\ell},\quad 0<\ell<1.
\]
The $x_i$ are the clockwise fiber generators in the algebraic convention
of \cite{NegamiKZ}. The $a_i$ are eigenangles of these group matrices,
not coefficients of logarithmic residues with an unspecified loop
orientation. Reversing a loop inverts its matrix. Write $\eta$ for the
unitary character on the chosen base subgroup. A different splitting
changes the recorded base words and character together.

The root action is pure, so it acts trivially on the abelianization of
$F_n$. Thus arbitrary scalar $q_i$ extend to this semidirect product;
they need not be equal. The KLM construction of \cite{HN} applies to this
Artin action. Its fiberwise middle-convolution realization, and base
change along the root map, identify the resulting local system with the
directional topological middle convolution of the seed. This uses the
bundle construction and compatibility of \cite{AH}; no unrestricted
claim about an additive Riemann--Hilbert comparison at resonant residue
parameters is needed.

\begin{remark}[The group is recorded explicitly]
The output is a representation of $\pi_1(M(\B))$. It factors through
$\pi_1(M(\A))$ exactly when every added-hyperplane meridian acts trivially.
If $\A=\A^Y$, no enlargement is required. For a reflection arrangement,
the complement group is the pure Artin group; extension to the full
Artin group requires additional equivariance. Naming the input Dynkin
type does not supply either descent or that equivariance.
\end{remark}

\subsection{The actual MC matrices and their middle quotient}
On $W=\C^n$, let $G_i$ be the identity except in row $i$, where
\begin{equation}\label{mc:eq:G}
 (G_i)_{ij}=\begin{cases}
 \lambda(q_j-1),&j<i,\\
 \lambda q_i,&j=i,\\
 q_j-1,&j>i.
 \end{cases}
\end{equation}
These are the scalar-input Dettweiler--Reiter/KLM convolution matrices
in the convention of \cite[Definition 9]{NegamiKZ}.
Define
\[
 K=\operatorname{span}\{e_i:q_i=1\},\qquad
 L=\bigcap_i\ker(G_i-I),\qquad V_{\MC}=W/(K+L).
\]
The matrices used as gates below are induced on this quotient.
Passing to $W/K$ deletes the indices with $q_i=1$; their fiber operators
become identity. Let $p=\#\{i:q_i\ne1\}$. Put
\[
 \delta=\begin{cases}1,&\lambda\prod_iq_i=1,\\0,&\text{otherwise.}\end{cases}
 \qquad \dim V_{\MC}=p-\delta.
\]
The case $p-\delta=0$ is a zero output, not a gate space.

To include the base group, a positive adjacent root interchange has
colored transport $T_i(q)$ equal to the identity away from indices
$i,i+1$, with block
\begin{equation}\label{mc:eq:T}
 T_i(q)\big|_{i,i+1}=\begin{pmatrix}0&q_i\\1&1-q_{i+1}\end{pmatrix}.
\end{equation}
It maps the fiber with colors $q$ to that with colors $s_iq$.
For a negative interchange use $T_i(s_iq)^{-1}$.
Along a chronological word the next matrix multiplies on the left,
and the colors are exchanged at each step. A pure word returns the
colors to their starting values. Multiply its resulting transport by
$\eta(\gamma)$ for a base loop $\gamma$. These formulas include the
base monodromy; using only the $G_i$ would specify only the fiber subgroup.
The actual base braid words must come from the given arrangement and
chosen based paths. The code accepts those words; it does not infer
all meridians of an arbitrary arrangement from unlabeled equations.

\subsection{Negami's invariant form, specialized to rank one}
Fix $s=\lambda^{1/2}=e^{\pi i\ell}$. Specializing the input form to $1$
in \cite[Theorems 8--9]{NegamiKZ} gives
\begin{equation}\label{mc:eq:H}
 H_{ij}=\begin{cases}
 s^{-1}(q_i^{-1}-\lambda)(q_i-1),&i=j,\\
 s^{-1}(q_i^{-1}-1)(q_j-1),&i<j,\\
 s(q_i^{-1}-1)(q_j-1),&i>j.
 \end{cases}
\end{equation}
This form descends to the middle quotient used below.

\begin{proposition}[Full monodromy invariance]\label{mc:prop:invariance}
The matrix $H$ is Hermitian, $G_i^*HG_i=H$, and
\begin{equation}\label{mc:eq:colored}
 T_i(q)^*H(s_iq)T_i(q)=H(q).
\end{equation}
It follows that the entire arrangement monodromy preserves $H$, including
the scalar base character. Its radical is $K+L$, and it induces a
nondegenerate form $H_{\MC}$ on $V_{\MC}$.
\end{proposition}
\begin{proof}
Hermitian symmetry follows from $\overline q_i=q_i^{-1}$ and
$\overline s=s^{-1}$. The fiber identity is the scalar case of
\cite[Theorem 8]{NegamiKZ}, or follows directly from
\eqref{mc:eq:G}--\eqref{mc:eq:H}. For unequal weights one must also verify
\eqref{mc:eq:colored}: outside the adjacent two rows and columns it is
unchanged; in the adjacent block it follows by multiplication of the
displayed $2\times2$ matrix, and the mixed rows give the exchanged
factors $q_i-1,q_{i+1}-1$. This identity telescopes along a pure word.
The unit-modulus scalar $\eta$ does not change the form.
For compatibility with the fiber action, writing $G'_j=G_j(s_iq)$,
the same calculation gives
\[
 T_iG_i=G'_{i+1}T_i,\quad
 T_iG_{i+1}=(G'_{i+1})^{-1}G'_iG'_{i+1}T_i,\quad
 T_iG_j=G'_jT_i\quad(j\ne i,i+1).
\]
These are the colored Artin-action identities; the colored braid
relations are those of the Long--Moody construction.
After deleting $K$, all $q_i\ne1$. If $B$ is the matrix whose $i$th
row is the nonzero row of $G_i-I$, then
\[
 H=\operatorname{diag}\bigl(s^{-1}(q_i^{-1}-1)\bigr)B.
\]
Its diagonal prefactor is invertible, so $\ker H=\ker B=L$ on $W/K$.
The radical statement on $W$ also follows from
\cite[Theorem 9]{NegamiKZ}. This proves the quotient assertion.
\end{proof}

\subsection{Exact signature and the gate condition}
Put
\[
 T=\ell+\sum_i a_i,\qquad m=\lfloor T\rfloor,
 \qquad \delta=\mathbf1_{T\in\mathbb Z}.
\]
\begin{theorem}[Rank-one specialization of Negami's signature formula]
For $\lambda\ne1$, the inertia of \eqref{mc:eq:H} and the signature of its
middle quotient are respectively
\begin{equation}\label{mc:eq:signature}
 (n_+,n_-,n_0)(H)=(p-m,\ m-\delta,\ n-p+\delta),
 \qquad \operatorname{sig}(H_{\MC})=(p-m,\ m-\delta).
\end{equation}
For a nonzero quotient, $H_{\MC}$ is positive definite exactly when
$m=\delta$, and negative definite exactly when $m=p$.
Equivalently, positive definiteness holds for $0<T\le1$, and negative
definiteness for $p\le T<p+1$, always with $p-\delta>0$.
\end{theorem}
\begin{proof}
In \cite[Theorem 1.1]{NegamiSignature}, the input dimension is $N=1$,
the eigenangle of $\lambda\prod_iq_i$ is $\{T\}$, the total kernel
dimension of the seed generators is $n-p$, and the product resonance
dimension is $\delta$. The integer in that formula becomes
$\ell+\sum_i a_i-\{T\}=\lfloor T\rfloor=m$, giving
\eqref{mc:eq:signature}. The definiteness assertions follow by setting the
appropriate negative or positive index to zero.
\end{proof}

For an elementary check, after deleting inactive indices, every initial
$k\times k$ principal minor is
\begin{equation}\label{mc:eq:minor}
 \det H_{[k]}=4^k\sin^{k-1}(\pi\ell)
 \sin\!\left(\pi\left(\ell+\sum_{j=1}^k a_j\right)\right)
 \prod_{j=1}^k\sin(\pi a_j).
\end{equation}
Indeed, writing $D=\operatorname{diag}(q_i-1)$ gives $H=D^*MD$,
where $M_{ii}=\sin\pi(a_i+\ell)/\sin\pi a_i$ and
$M_{ij}=s^{-1}$ for $i<j$, $s$ for $i>j$. Subtracting successive
rows after clearing the diagonal denominators, or induction by a Schur
complement, gives \eqref{mc:eq:minor}. It is an additional check on the
normalization; no nonsingular pivot is assumed at resonance.

\begin{proposition}[Necessity for the complete output]
For a nonzero rank-one MC output, its fiber representation is irreducible.
Consequently the full arrangement representation is unitarizable if and
only if $H_{\MC}$ is definite (with either sign).
\end{proposition}
\begin{proof}
Work after deleting inactive indices. Let $\overline W$ be a nonzero
invariant subspace of $\C^p/L$, and lift a nonzero class to $v\notin L$.
Some $(G_i-I)v$ is a nonzero multiple of $e_i$. The class of $e_i$ is
nonzero: for $p\ge2$ it is not fixed by $G_j$, $j\ne i$, since
$q_i-1\ne0$; for $p=1$ the nonzero quotient is one dimensional.
For every $j\ne i$, $(G_j-I)e_i$ is a nonzero multiple of $e_j$.
Thus $\overline W$ contains all coordinate classes and is the full
quotient. Any invariant positive form, compared with the nondegenerate
$H_{\MC}$, gives a commuting endomorphism; Schur's lemma makes that
endomorphism scalar. Hence a positive invariant form exists precisely
when $H_{\MC}$ is definite.
\end{proof}

Choose an orthonormal matrix $Q$ whose columns span $(K+L)^\perp$.
The induced matrices and form are
\[
 \widehat M_\gamma=Q^*M_\gamma Q,\qquad H_Q=Q^*HQ.
\]
The compression is the quotient action, not a claim that this orthogonal
complement is invariant. Choose $\varepsilon=\pm1$ so that
$\varepsilon H_Q>0$, and choose $S$ with $S^*S=\varepsilon H_Q$.
Then the actual MC gates are
\begin{equation}\label{mc:eq:gates}
 U_\gamma=S\widehat M_\gamma S^{-1},\qquad U_\gamma^*U_\gamma=I.
\end{equation}
All generators use this one common basis. A $D$-dimensional output is
a $D$-level gate system; it is directly a $k$-qubit system only when
$D=2^k$, unless an explicit larger-space encoding is supplied.

\subsection{Non-Clifford gates, closure, and finite image}
For an active fiber index, the nontrivial eigenvalue is
$\mu_i=\lambda q_i$. On a definite output of dimension at least two,
\[
 U_i=I+(\mu_i-1)v_iv_i^*,\qquad \|v_i\|=1.
\]
Here $v_i$ is the normalized vector $SQ^*e_i$.
The inner products of these directions are, up to positive normalization,
the entries of $\varepsilon H$. In particular all pairwise inner products
are nonzero, and the directions span the output space.

\begin{theorem}[A sufficient full-closure condition]\label{mc:thm:dense}
Suppose the quotient form is definite, its dimension is $D\ge2$, and
$\ell+a_i$ is irrational for every active index. Then the closure of
the fiber image, and therefore of the full arrangement image after
\eqref{mc:eq:gates}, is $\U(D)$. Every active fiber gate has infinite
projective order and is non-Clifford in every finite Weyl--Heisenberg frame.
\end{theorem}
\begin{proof}
The powers of $U_i$ are dense in the circle
$\{\exp(itv_iv_i^*):t\in\mathbb R\}$. Thus the closure Lie algebra
contains every $iv_iv_i^*$. Rank-one projectors onto two independent,
nonorthogonal directions generate $\mathfrak u(2)$ on their span:
their commutator and a second commutator give the two off-diagonal
directions, and their traces give the central direction. Inductively,
if $\mathfrak u(W)$ is available and a new normal has both a nonzero
component in $W$ and outside $W$, the same commutators add
$\mathfrak u(W+\C v)$. Pairwise nonzero inner products and spanning
therefore give $\mathfrak u(D)$. The closed subgroup is $\U(D)$.
The eigenvalue ratio $\mu_i$ is not a root of unity. The projective
normalizer of any finite irreducible Weyl--Heisenberg group is finite:
it acts by automorphisms on that finite group, and its kernel is scalar by
irreducibility. Hence an infinite-projective-order gate cannot be Clifford.
This uses the same spectral principle as \cite{NegamiGates}.
\end{proof}

For example, choose all active $a_i$ rational and positive with
$\sum_i a_i<1$, and choose irrational $0<\ell<1-\sum_i a_i$.
Then the quotient has dimension $p$, is positive definite, and the
theorem applies. It works for every input arrangement and direction
with $p\ge2$, on its stated directional closure; this assertion does
not require surjectivity onto a universal configuration-space group.

\begin{theorem}[Exact finite-image test for rational seed phases]
Assume all $a_i,\ell$ are rational, $\lambda\ne1$, the base character
$\eta$ has finite image, and $D=p-\delta>0$. Choose a common denominator
$N$. For every $k$ coprime to $N$, define
\[
 T_k=\{k\ell\}+\sum_i\{ka_i\},\qquad m_k=\lfloor T_k\rfloor.
\]
The full arrangement MC image is finite if and only if
\begin{equation}\label{mc:eq:finite}
 m_k\in\{\delta,p\}\quad\text{for every }k\in(\mathbb Z/N\mathbb Z)^\times.
\end{equation}
\end{theorem}
\begin{proof}
Under a cyclotomic embedding the scalar eigenangles become their
fractional multiples, while inactive indices and the resonance dimension
are unchanged. Formula \eqref{mc:eq:signature} gives the conjugate quotient
signature $(p-m_k,m_k-\delta)$.
If the arrangement image is finite, its fiber image is finite. Its
conjugate fiber representations are irreducible by the preceding
proposition. Averaging over each finite image and applying Schur's lemma
forces the canonical conjugate form to be definite. This proves necessity.
Conversely, the ambient matrices \eqref{mc:eq:G}, the colored transports,
and their inverses have cyclotomic integral entries, since the scalar
parameters are roots of unity. After extending the field to include
the finite values of $\eta$, their quotient preserves the image of a
full integral lattice. The quotient is defined over that field, including
at resonance. If \eqref{mc:eq:finite} holds, every complex embedding preserves
a positive form after a possible sign change. The diagonal embedding
of the lattice in the product of the archimedean spaces is discrete,
whereas the matrix images lie in a compact product of unitary groups.
Only finitely many images of a lattice basis are possible. This proves
finiteness. This is the scalar specialization of the arithmetic approach
in \cite{NegamiFinite}, with the argument included to specify its scope.
\end{proof}

The exact test does not infer finite image from finite orders of the
individual generators. For example, $n=3$, all $a_i=\ell=1/10$, has
positive form at the original embedding, but $k=3$ gives
$T_k=6/5$, hence signature $(2,1)$: the image is infinite.
The cases $n=2$, all $a_i=\ell=1/6$, and $n=3$, all
$a_i=\ell=1/4$ (a resonant rank-two quotient), pass the finite test.
Theorem~\ref{mc:thm:dense} is a sufficient closure theorem. We do not
classify all closures for rational parameters or for mixtures of
rational and irrational eigenangles. At rational parameters, Clifford
membership also requires a specified computational frame; the code's
Weyl-normalizer distances are numerical diagnostics, not exact certificates.

\subsection{An arrangement example and Dynkin inputs}
Start with the genuinely affine three-line arrangement
\[
 \A=\{y=0,\ y=x,\ y=1-x\},\qquad
 \B=\A\cup\{x=0,x=1,x=\tfrac12\}.
\]
Use the unitary rank-one seed with fiber eigenangles
\[
 (a_1,a_2,a_3)=(1/20,1/15,1/12),\quad
 \ell=\sqrt2/20,
\]
and trivial character on the base subgroup. The fiber weight sum is
$1/5$, so $T=1/5+\sqrt2/20<1$. There is no middle rank drop, and
the actual three-dimensional MC output is positive definite and has
closure $\U(3)$. Each of its three fiber gates is non-Clifford by
Theorem~\ref{mc:thm:dense}. These are qutrit gates.
For this example the fiber matrices used by the program are explicitly
\[
 G_1=\begin{pmatrix}\lambda q_1&q_2-1&q_3-1\\0&1&0\\0&0&1\end{pmatrix},\quad
 G_2=\begin{pmatrix}1&0&0\\\lambda(q_1-1)&\lambda q_2&q_3-1\\0&0&1\end{pmatrix},
\]
\[
 G_3=\begin{pmatrix}1&0&0\\0&1&0\\\lambda(q_1-1)&\lambda(q_2-1)&\lambda q_3\end{pmatrix}.
\]
The root map of the base is $x\mapsto(0,x,1-x)$. Base meridians
are evaluated by their based pure braid words in \eqref{mc:eq:T}; their
precise stems affect the conjugating words, but not the common
invariant form or the full-closure conclusion, which already follows
from the actual fiber subgroup.

Dynkin types can specify the input arrangement without replacing this
MC construction. For a generic seed with all $q_i\ne1$ and no product
resonance, its output rank is the number $n$ of transverse hyperplanes,
not the rank of the root system:
\begin{center}
\begin{tabular}{llll}
\toprule
Input & Direction & $n$ & Enlargement\\
\midrule
$A_r$ & forget the last labelled coordinate & $r$ & none\\
$B_r=C_r$ & last coordinate & $2r-1$ & none\\
$D_r$ & last coordinate & $2r-2$ & add $x_i=0$, $i<r$\\
$E_6$ & generic direction & $36$ & collision closure\\
$E_7$ & generic direction & $63$ & collision closure\\
$E_8$ & generic direction & $120$ & collision closure\\
\bottomrule
\end{tabular}
\end{center}
For $B_r$, the punctures are $0,\pm x_1,\ldots,\pm x_{r-1}$,
and all their collisions are already base hyperplanes of type $B_{r-1}$.
For $D_r$ the punctures are $\pm x_1,\ldots,\pm x_{r-1}$;
the collisions $x_i=-x_i$ add the missing coordinate hyperplanes.
A generic direction for $E_r$ is transverse to every root hyperplane,
giving the displayed counts. Special directions can change $n$.
In particular, this construction does not turn an $E_8$ rank-one seed
into eight-dimensional gates merely because the root system has rank eight.

The Burau terminology refers to the rank-one homological construction:
on the universal configuration bundle the corresponding colored
representation is the reduced Gassner-type twisted homology representation.
When the $n$ seed colors and the convolution color are all equal,
its nonresonant specialization recovers reduced Burau on the pure braid
group, with the appropriate loop and basis conventions \cite{HN}.
Pulling back to an arrangement gives the Burau-type family discussed here.
At resonance the middle quotient must be taken; the full homological
Burau specialization and that quotient need not have equal dimensions.
No equality with Squier's general Coxeter-type construction is assumed
or needed for the gates in this note.

\subsection{Reproducibility of the MC gate calculations}
Run \texttt{python3 mc\_arrangement\_gates.py --verify --verify-symbolic}.
The script builds the MC matrices, the form, its radical quotient and
one common unitarization, and exports a qutrit gate example. The symbolic
checks use rational Laurent expressions in independent colors and the
square root of $\lambda$; numerical tests check 164 tuples, including
zero seed angles, resonance, and negative-definite cases. Rational
finite-image tests use exact fractions and modular enumeration.
These checks supplement the stated proofs; they do not replace their
review. The group closure arguments are mathematical arguments, not
inferences from numerical word searches. Large group enumeration is
not required.

This working note and its code were drafted with substantial assistance
from OpenAI ChatGPT. It explicitly applies the author's prior invariant
form and signature results. The recorded verification concerns only the
checks listed above and does not certify independent human review or
the novelty of the results. Sections~1--6 use the Coxeter--Burau construction; the present
section uses the MC matrices \eqref{mc:eq:G} on their middle quotient.


\appendix
\section{Writing reflection arrangements by Dynkin type}\label{sec:dynkin}
\subsection*{The scope of the type notation}
A general affine arrangement is specified by equations
\[
\A=\{H_a:\ell_a(x)=b_a\}_{a=1}^N,
\qquad Q_{\A}(x)=\prod_{a=1}^N(\ell_a(x)-b_a).
\]
There is no rule assigning a Dynkin type to every such arrangement.
The notation ``type $X_r$'' used here means that, after a linear change
of coordinates and removal of any redundant ambient factor, the
arrangement is the complexified reflection arrangement of a root system
$\Phi(X_r)$.  A common translation is allowed if the original equations
are affine.  In particular, counting hyperplanes or drawing their
intersection graph does not determine a Dynkin type.

For a reduced crystallographic root system $\Phi\subset V_{\mathbb R}^*$,
the construction is
\begin{equation}\label{eq:root-arrangement}
\boxed{\quad \A(X_r)=\{\ker\alpha:\alpha\in\Phi^+(X_r)\},
\qquad Q_{X_r}(x)=\prod_{\alpha\in\Phi^+(X_r)}\alpha(x).\quad}
\end{equation}
Here $V=V_{\mathbb R}\otimes\C$, and one positive root is used from
each pair $\{\alpha,-\alpha\}$.  Thus the Dynkin diagram has $r$
vertices, whereas the arrangement has $|\Phi|/2$ hyperplanes.
The root-system construction is implemented in \cite{SageArr}.

The irreducible crystallographic types are
$A_n,B_n,C_n,D_n,E_6,E_7,E_8,F_4,G_2$; reducible types give products
of arrangements.  See \cite[Theorem 23.7]{Etingof} for the classification.
Finite real reflection arrangements also include the noncrystallographic
Coxeter types $H_3,H_4,I_2(m)$, for which one uses Coxeter rather than
crystallographic Dynkin data.

\subsection*{Classical types in coordinates}
For $A_n$, use
$V=\{(x_1,\ldots,x_{n+1})\in\C^{n+1}:\sum_i x_i=0\}$.
For $B_n,C_n,D_n$, use $V=\C^n$.  In the table, every displayed linear
form is set equal to zero, with $i<j$ in the appropriate coordinate range.
\begin{center}
\renewcommand{\arraystretch}{1.2}
\begin{tabular}{lll}
\toprule
Type & Defining linear forms & Number of hyperplanes\\
\midrule
$A_n$ & $x_i-x_j$ & $n(n+1)/2$\\
$B_n$ & $x_i,\ x_i-x_j,\ x_i+x_j$ & $n^2$\\
$C_n$ & $2x_i,\ x_i-x_j,\ x_i+x_j$ & $n^2$\\
$D_n$ & $x_i-x_j,\ x_i+x_j$ & $n(n-1)$\\
\bottomrule
\end{tabular}
\end{center}
These follow by taking kernels of the standard classical roots
\cite[\S20.3]{Etingof}; the $B_n$ equations are also listed explicitly
in \cite{M2B}.  Consequently, up to nonzero constants,
\begin{align*}
Q_{A_n}&=\prod_{1\le i<j\le n+1}(x_i-x_j),\\
Q_{B_n}=Q_{C_n}&=\left(\prod_{i=1}^n x_i\right)
                 \prod_{1\le i<j\le n}(x_i^2-x_j^2),\\
Q_{D_n}&=\prod_{1\le i<j\le n}(x_i^2-x_j^2).
\end{align*}
The $A_n$ polynomial is restricted to $V$ above; using all of
$\C^{n+1}$ instead gives the familiar nonessential braid arrangement
with a redundant diagonal direction.

For example,
\[
\A(B_3)=\{x_1=0,x_2=0,x_3=0\}
\cup\{x_i-x_j=0,x_i+x_j=0:1\le i<j\le3\}
\]
has nine planes, even though its Dynkin diagram has three vertices.
The arrangement of type $D_3$ has six planes and is isomorphic to
type $A_3$; the irreducible $D_n$ family is usually indexed by $n\ge4$.

\paragraph{Why $B_n$ and $C_n$ give the same arrangement.}
Their root systems use $\pm e_i$ and $\pm2e_i$, respectively, together
with $\pm e_i\pm e_j$.  Multiplying a normal by $2$ does not change
its kernel.  Hence
\[
\A(B_n)=\A(C_n)
\]
in these coordinates.  The hyperplane set alone therefore does not
retain the distinction between these dual root systems or the direction
of the double-edge arrow in their Dynkin diagrams.

\subsection*{Recovering the $E$-type equations from the diagram}
One need not list coordinates in $\mathbb R^8$ to specify the exceptional
arrangements.  Choose simple roots $\alpha_1,\ldots,\alpha_r$, and use
the coordinates $z_i=\alpha_i(x)$.  If a positive root has simple-root
coefficients $c=(c_1,\ldots,c_r)$, its hyperplane is exactly
\begin{equation}\label{eq:root-coefficient-hyperplane}
H_c=\{c_1z_1+\cdots+c_rz_r=0\}.
\end{equation}

For $E_r$, $r=6,7,8$, use the diagram with edge set
\[
\mathcal E_r=\{\{i,i+1\}:1\le i\le r-2\}\cup\{\{3,r\}\}.
\]
For $r=6$ this agrees with the numbering used in
Section~\ref{sec:general-burau}.  Its Cartan matrix is
\[
C_{ii}=2,\qquad C_{ij}=-1\ \text{if }\{i,j\}\in\mathcal E_r,
\qquad C_{ij}=0\ \text{otherwise}.
\]
Start with the unit vectors $e_i$ in the simple-root coefficient lattice
and repeatedly apply
\begin{equation}\label{eq:root-generation}
s_i(c)=c-\left(\sum_j C_{ij}c_j\right)e_i.
\end{equation}
Closing under these reflections gives all roots; retain those whose
coefficients are all nonnegative, and apply
\eqref{eq:root-coefficient-hyperplane}.  This is the usual simple-reflection
generation of a finite root system \cite[\S\S21--23]{Etingof}, written
as an explicit equation-generation procedure.

The following counts distinguish arrangement size from rank:
\begin{center}
\begin{tabular}{lrrr}
\toprule
Type & Rank & Roots & Hyperplanes\\
\midrule
$E_6$ & $6$ & $72$ & $36$\\
$E_7$ & $7$ & $126$ & $63$\\
$E_8$ & $8$ & $240$ & $120$\\
$F_4$ & $4$ & $48$ & $24$\\
$G_2$ & $2$ & $12$ & $6$\\
\bottomrule
\end{tabular}
\end{center}
The three $E$-type counts were also checked in the original note by the finite orbit procedure
\eqref{eq:root-generation}.  The same root-kernel recipe
\eqref{eq:root-arrangement} applies to $F_4$ and $G_2$.

\subsection*{From the Dynkin diagram to the Burau matrices}
For the representation in Section~\ref{sec:general-burau}, the relevant
data are the Coxeter labels.  The translation is
\begin{center}
\begin{tabular}{lcc}
\toprule
Dynkin adjacency & $m_{ij}$ & Coefficient of $e_i$ in $R_i e_j$\\
\midrule
No edge & $2$ & $0$\\
Single edge & $3$ & $t$\\
Double edge & $4$ & $\sqrt2\,t$\\
Triple edge & $6$ & $\sqrt3\,t$\\
\bottomrule
\end{tabular}
\end{center}
This is direct substitution into \eqref{eq:coxeter-burau} for $i\ne j$.
For $B_n$ and $C_n$, the Coxeter graph is the path with
$m_{i,i+1}=3$ for $i<n-1$ and $m_{n-1,n}=4$.
Thus both types give the same matrices in this symmetric normalization.
The arrow records root lengths, which this Coxeter-matrix construction
does not require.

To construct a crystallographic arrangement in SageMath \cite{SageArr}:
\begin{verbatim}
A = hyperplane_arrangements.Coxeter("B3")
A.hyperplanes()
A.defining_polynomial()
E = hyperplane_arrangements.Coxeter("E6")
\end{verbatim}
The constructor may use a redundant ambient space (for example for
$A_n$ or $E_6$); its essential rank is the Dynkin subscript.

\paragraph{Notation and limits.}
The ordinary braid group on $n$ strands is of type $A_{n-1}$:
$\operatorname{Art}(A_{n-1})\cong\mathsf B_n$.
It should not be confused with $\operatorname{Art}(B_n)$, the Artin
group of Dynkin type $B_n$.
Also, a general affine arrangement is not automatically an affine
Dynkin arrangement.  For instance, the triangle
$\{y=0,y=x,y=1-x\}$ used later has no common intersection, so it
cannot be translated into a finite central reflection arrangement.
Its equations remain valid input data without a Dynkin label.

\section{Scope and existing constructions}\label{sec:general-arrangement}
Let $\A$ be a finite arrangement of affine hyperplanes in $\C^\ell$, and put
$M(\A)=\C^\ell\setminus\bigcup_{H\in\A}H$.
In these appendices $\kappa$ denotes the additive convolution parameter,
not the Coxeter gate eigenvalue in Sections~1--6 or the
multiplicative convolution parameter in Section~\ref{sec:mc-gates}.
Fix a one-dimensional direction $Y$.  The minimal directional closure
$\B=\A^Y$ adds the collision hyperplanes necessary to make the projection
along $Y$ a bundle of punctured affine lines.  This closure and the
corresponding middle convolution are part of the existing theory
\cite{AH,Hiroe}.  The relationship between Long--Moody induction and Katz
middle convolution, including applications to fiber-type arrangements,
was developed in \cite{HN}.  Earlier homological generalizations of
Magnus, Burau, and Gassner representations occur in \cite{CS}.

Two distinctions will be maintained throughout.  First, the full convolution
space is not automatically the middle-convolution space: the invariant
subspace $K+L_\kappa$ must be divided out.  Second, a representation on
$M(\B)$ does not automatically descend to $M(\A)$.  Neither passing to the
middle quotient nor writing down all meridian matrices removes the second
issue.

\section{The algebraic construction in arbitrary rank}
The holonomy Lie algebra has the presentation
\begin{equation}\label{eq:holonomy}
\h(\A)=\operatorname{Lie}\langle t_H:H\in\A\rangle\big/
\left\langle\left[t_H,\sum_{H'\supset X}t_{H'}\right]:
 X\in L_2(\A),\ H\supset X\right\rangle .
\end{equation}
Here $L_2(\A)$ consists of the nonempty intersections of codimension two.
Thus an input is a finite-dimensional vector space $E$ and endomorphisms
$A_H$ satisfying the displayed commutator relations.  Extend the input to
$\B$ by setting $A_D=0$ for $D\in\B\setminus\A$.

Choose coordinates $(x,y)$ with $Y$ the $y$-direction.  Write the transverse
hyperplanes as $H_i=\{y=f_i(x)\}$, $1\leq i\leq n$; all other hyperplanes
are parallel to $Y$.  The functions $f_i$ are affine linear.  A nonconstant
difference $f_i-f_j$ defines a collision hyperplane in the base, whose
inverse image is included in $\B$.  Constant nonzero differences give no
collision divisor.  We identify a base hyperplane with its inverse image
when indexing residues.

Set $W=E^n$ and use blocks indexed by $i$, with the coordinates of $E$
inside each block.  The known convolution formulas specialize to
\begin{align}
C_{H_i}&=\sum_{j=1}^n E_{ij}\otimes A_{H_j}
          +\kappa E_{ii}\otimes\Id_E,\label{eq:transverse}\\
C_D&=I_n\otimes A_D+
\sum_{\substack{i<j\\\{f_i=f_j\}=D}}
\bigl((E_{ii}-E_{ij})\otimes A_{H_j}
 +(E_{jj}-E_{ji})\otimes A_{H_i}\bigr).
\label{eq:parallel}
\end{align}
These are the block form of \cite[Proposition~3.1]{AH}.  They define a representation
of $\h(\B)$.  The compatibility with the Lie-algebraic convolution functor
is \cite[Theorem 4.7]{Hiroe}; constructing these matrices requires no
integration of a differential equation.

The two invariant subspaces and the middle quotient are
\begin{equation}\label{eq:middle}
K=\bigoplus_i e_i\otimes\ker A_{H_i},\qquad
L_\kappa=\ker\Bigl(\sum_i C_{H_i}\Bigr),\qquad
\mc^Y_\kappa(E)=W/(K+L_\kappa).
\end{equation}
For $\kappa\ne0$,
\[
L_\kappa=\{(v,\ldots,v):(\sum_i A_{H_i}+\kappa\Id_E)v=0\},
\qquad K\cap L_\kappa=0.
\]
In particular, if $r=\dim E$, the output dimension is
\[
nr-\sum_i\dim\ker A_{H_i}
 -\dim\ker(\sum_i A_{H_i}+\kappa\Id_E).
\]
Constant intertwiners induce intertwiners on the quotient, giving
\begin{equation}\label{eq:functor}
\mc^Y_\kappa:\Rep\h(\A)\longrightarrow\Rep\h(\A^Y).
\end{equation}
This is an endofunctor on the original representation category when
$\A=\A^Y$.  Otherwise, the codomain in \eqref{eq:functor} must not be
silently replaced by $\Rep\h(\A)$.

\section{Equal weights and the Burau pullback}\label{sec:equal}
Take $E=\C$, $A_{H_i}=\alpha$, $A_D=0$ on every original parallel
hyperplane, and $\kappa=\alpha\ne0$.  Then
\begin{equation}\label{eq:equal}
C_{H_i}=\alpha e_i(\one^{\mathsf T}+e_i^{\mathsf T}),\qquad
C_D=\alpha\sum_{\substack{i<j\\\{f_i=f_j\}=D}}
 (e_i-e_j)(e_i-e_j)^{\mathsf T}.
\end{equation}
The quotient in \eqref{eq:middle} is still explicitly checked:
$K=0$, and $L_\alpha=0$ because $(n+1)\alpha\ne0$.  Hence the
\emph{additive} middle-convolution space has dimension $n$ in this case.

Let $m=n+1$ and define
\begin{equation}\label{eq:F}
F:M(\B)\longrightarrow\Conf_m(\C),\qquad
F(x,y)=(f_1(x),\ldots,f_n(x),y).
\end{equation}
All points on the right are distinct by construction of $\B$.
On $V=\C^m/\C\one$, use quotient coordinates $w_i-w_m$.
Let $T_{ij}$ be the induced action of
$(e_i-e_j)(e_i-e_j)^{\mathsf T}$ from $\C^m$ on $V$.

\begin{proposition}\label{prop:gauge}
Set $S=I_n+J_n$, where $J_n$ is the all-ones matrix.  In the above
quotient coordinates,
\[
S(C_{H_i}/\alpha)S^{-1}=T_{i,m},\qquad
S\bigl((e_i-e_j)(e_i-e_j)^{\mathsf T}\bigr)S^{-1}=T_{ij}
\quad(i,j\le n).
\]
Consequently, the connection $d-\sum_H C_H\,d\log f_H$ is, after
this constant change of basis, the pullback of
$d-\alpha\sum_{i<j}T_{ij}\,d\log(z_i-z_j)$.
\end{proposition}
\begin{proof}
In quotient coordinates, $T_{i,m}=(\one+e_i)e_i^{\mathsf T}$.
Since $S e_i=\one+e_i$ and $e_i^{\mathsf T}S=\one^{\mathsf T}+e_i^{\mathsf T}$,
the first identity follows.  Every $e_i-e_j$ has coordinate sum zero,
so $S$ acts as the identity on it and commutes with its outer product.
The logarithmic forms pull back to $d\log(y-f_i)$ and
$d\log(f_i-f_j)$, respectively.  Grouping proportional differences
gives \eqref{eq:parallel}.
\end{proof}

Put $q=\exp(2\pi i\alpha)$.  For generic parameters, this is the
homological reduced Burau representation restricted to $P_m$ and pulled
back by $F_*$.  One sufficient regime for the period comparison used here
is $\alpha\notin\mathbb Z$ and $m\alpha\notin\mathbb Z$.
The relevant local system has common monodromy $q$ at the $m$ finite
punctures; its twisted first homology has dimension $m-1$.
The connection/homology comparison is part of the standard framework
reviewed in \cite[\S3]{Kohno}, together with the convolution comparison
of \cite[Theorem 4.2]{AH}.

\begin{remark}
At resonant parameters, one must distinguish the full homological Burau
module, its algebraic specialization, and the middle-convolution quotient.
For example, trivial monodromy at infinity can contribute a boundary
subspace to the topological middle quotient.  The generic identification
above is not a claim that all these spaces have the same rank at every
special parameter.
\end{remark}

\section{Closed formulas for based pure-braid matrices}
Use the unreduced Burau convention with block
$\left(\begin{smallmatrix}1-q&q\\1&0\end{smallmatrix}\right)$ for
$\sigma_i$, and then pass to $\C^m/\C\one$.  We act on column vectors;
in products of paths the right-hand path is traversed first.  Positive
half twists are counterclockwise.  Write
\[
A_{ij}=\sigma_{j-1}\cdots\sigma_{i+1}\sigma_i^2
       \sigma_{i+1}^{-1}\cdots\sigma_{j-1}^{-1}.
\]

\begin{proposition}\label{prop:pure}
Let $n=m-1$ and $1\le i<j\le m$.  In quotient coordinates,
\begin{equation}\label{eq:pure}
\beta_q(A_{ij})=I_n+(q-1)u_{ij}v_{ij}^{\mathsf T}.
\end{equation}
For $j\le n$,
\[
u_{ij}=q e_i+(q-1)\sum_{i<k<j}e_k-e_j,
\qquad v_{ij}=e_i-e_j.
\]
For $j=m$,
\[
u_{i,m}=\sum_{k<i}e_k+(q+1)e_i+q\sum_{k>i}e_k,
\qquad v_{i,m}=e_i.
\]
In particular,
\[
\det\beta_q(A_{ij})=q^2,\qquad
(\beta_q(A_{ij})-I)(\beta_q(A_{ij})-q^2I)=0.
\]
\end{proposition}
\begin{proof}
For $i<n$, the square of the adjacent Burau block differs from $I$ by
$(q-1)(q e_i-e_{i+1})(e_i-e_{i+1})^{\mathsf T}$.
Conjugate this rank-one factorization by
$\beta_q(\sigma_{j-1}\cdots\sigma_{i+1})$.
Moving the last nonzero entry successively to the right produces the
intermediate coefficients $q-1$.  At the last strand, the quotient
generator has final column $(-1,\ldots,-1,-q)^{\mathsf T}$; applying it
gives the stated formula for $u_{i,m}$ and $v_{i,m}$.
The case $i=n,j=m$ follows by squaring this quotient generator directly.
Finally $v_{ij}^{\mathsf T}u_{ij}=q+1$, giving the determinant and
quadratic identities.
\end{proof}

These are based group matrices, not individual exponentials of residues.
For a general arrangement loop $\gamma$, one computes
$\rho(\gamma)=\beta_q(F_*\gamma)$ from its braid word.  Determining that
word includes the paths from a common basepoint.  At a multiple collision,
the local braid is a full twist on the colliding cluster; different
clusters and the connecting path must also be accounted for.

\section{A six-line example with fixed meridians}\label{sec:six-lines}
Start with $\A=\{y=0,y=x,y=1-x\}$ and project along $y$.
Then
\[
\B=\{y=0,y=x,y=1-x,x=0,x=1,x=\tfrac12\}.
\]
Choose basepoint $(x,y)=(1/4,2)$, labelled configuration
$(0,x,1-x,y)$, and positive local circles.  For the three transverse
meridians keep $x=1/4$ and take the $y$-stems below the real axis to just
right of the relevant puncture.  For the other meridians keep $y=2$:
the stems to $x=0$ and $x=1/2$ are real, ending respectively to the right
and left; the stem to $x=1$ passes above $x=1/2$ and ends to the left of
$1$.  Return along the reverse of each stem after the local circle.

With $d=q-1$, all six matrices are specified by
$M_H=I_3+d u_Hv_H^{\mathsf T}$ and the following table.
\begin{center}
\renewcommand{\arraystretch}{1.2}
\begin{tabular}{llll}
\toprule
$H$ & $F_*(m_H)$ & $u_H^{\mathsf T}$ & $v_H^{\mathsf T}$\\
\midrule
$y=0$ & $A_{14}$ & $(q+1,q,q)$ & $(1,0,0)$\\
$y=x$ & $A_{24}$ & $(1,q+1,q)$ & $(0,1,0)$\\
$y=1-x$ & $A_{34}$ & $(1,1,q+1)$ & $(0,0,1)$\\
$x=0$ & $A_{12}$ & $(q,-1,0)$ & $(1,-1,0)$\\
$x=1$ & $A_{13}$ & $(q,q-1,-1)$ & $(1,0,-1)$\\
$x=1/2$ & $A_{23}$ & $(0,q,-1)$ & $(0,1,-1)$\\
\bottomrule
\end{tabular}
\end{center}
Explicitly, the original hyperplanes have matrices
\begin{align*}
M_{y=0}&=\begin{pmatrix}q^2&0&0\\qd&1&0\\qd&0&1\end{pmatrix},&
M_{y=x}&=\begin{pmatrix}1&d&0\\0&q^2&0\\0&qd&1\end{pmatrix},&
M_{y=1-x}&=\begin{pmatrix}1&0&d\\0&1&d\\0&0&q^2\end{pmatrix}.
\end{align*}
The additional hyperplanes have matrices
\begin{align*}
M_{x=0}&=\begin{pmatrix}1+qd&-qd&0\\-d&q&0\\0&0&1\end{pmatrix},&
M_{x=1}&=\begin{pmatrix}1+qd&0&-qd\\d^2&1&-d^2\\-d&0&q\end{pmatrix},\\
M_{x=1/2}&=\begin{pmatrix}1&0&0\\0&1+qd&-qd\\0&-d&q\end{pmatrix}.&&
\end{align*}
These extra matrices are nonidentity for $q\ne1$, including $q=-1$.
Thus this particular group representation does not descend to the
original three-line complement for $q\ne1$.

As consistency checks, setting $S=I_3+J_3$ and denoting by $\widehat C_H$
the residue in \eqref{eq:equal} with $\alpha=1$, one obtains
\[
\left.\frac{dM_H}{dq}\right|_{q=1}=S\widehat C_HS^{-1}
\quad\text{for all six hyperplanes}.
\]
The standard pure-braid full-twist product also gives
\[
M_{12}M_{13}M_{23}M_{14}M_{24}M_{34}=q^4I_3.
\]
The first identity is a first-order check, not an identification of
$M_H$ with the exponential of its residue at finite parameter.

\section{The largest descended quotient and its obstruction}\label{sec:descent}
The natural surjection $p:\h(\B)\to\h(\A)$ kills precisely the Lie
ideal generated by $t_D$, $D\in\B\setminus\A$.  For a finite-dimensional
$\h(\B)$-module $V$ with action matrices $B_H$, define
\begin{equation}\label{eq:saturation}
N_0=\sum_{D\in\B\setminus\A}\im B_D,\qquad
N_{k+1}=N_k+\sum_{H\in\B}B_HN_k.
\end{equation}

\begin{proposition}\label{prop:descent}
The sequence in \eqref{eq:saturation} stabilizes after finitely many
strict dimension increases.  If its stable value is $N$, then $V/N$ is
the largest quotient of $V$ that factors through $\h(\A)$.
\end{proposition}
\begin{proof}
The stable space is the smallest invariant subspace containing all
$\im B_D$.  Thus every additional generator acts as zero on $V/N$, and
so does the ideal they generate.  Conversely, the kernel of any quotient
map on which those generators act trivially must contain this invariant
subspace.  Finite dimensionality gives termination.
\end{proof}

\begin{lemma}\label{lem:irred}
For $\alpha\ne0$, the operators
$C_i=\alpha e_i(\one^{\mathsf T}+e_i^{\mathsf T})$ on $\C^n$
act irreducibly.
\end{lemma}
\begin{proof}
The matrix whose $i$th row is $\one^{\mathsf T}+e_i^{\mathsf T}$ is
$I_n+J_n$, which is invertible.  Given a nonzero vector in an invariant
subspace, some $C_i$ therefore sends it to a nonzero multiple of $e_i$.
For $j\ne i$, $C_j e_i=\alpha e_j$.  Hence the subspace contains every
coordinate vector.
\end{proof}

\begin{corollary}\label{cor:zero}
For the equal-weight specialization of Section~\ref{sec:equal}, if $\A^Y\ne\A$,
then its largest quotient descending to $\h(\A)$ is zero.
\end{corollary}
\begin{proof}
An added hyperplane has at least one collision class of size at least
two.  Its residue is a nonzero multiple of the sum of the complete-graph
Laplacians on its disjoint collision classes, so it is nonzero.  The
invariant subspace generated by its image is nonzero and hence, by
Lemma~\ref{lem:irred}, equals the whole module.
\end{proof}

This corollary concerns Lie-algebra descent.  For an actual group
representation $\rho:G_{\B}\to\mathrm{GL}(V)$, the corresponding
largest descended quotient is
\[
V\big/\operatorname{span}\{\rho(g)(\rho(m_D)-I)v:
g\in G_{\B},\ D\in\B\setminus\A,\ v\in V\}.
\]
It kills additional meridians rather than additional residues.  In
particular, at resonant parameters the two criteria must not be
identified without an additional argument.


\section{Reproducibility of the arrangement calculations}
The accompanying \texttt{legacy\_arrangement/} directory contains the
Python/SymPy sources and regenerated reports.  Those reports record checks of
the block relations, middle quotients, invariant saturations, Fox
derivatives, 34 pure generators for $m=3,4,5,6$, the six first derivatives,
and the full-twist identity.  All four arrangement verification scripts
were rerun successfully for this public note package.
The path-tracking check of the six braid words uses floating-point
arithmetic and is not an interval-arithmetic certificate.
The separate Coxeter verification files are described in the main text.
Neither a finite word search nor floating-point matrix checks are used
as proofs of density or of infinitude.

The reflection arrangements of types $E_6,E_7,E_8$ are not fiber-type,
by the supersolvability classification \cite[Theorem~1.2]{HR}.
Thus results restricted to fiber-type arrangements do not alone cover
these cases.  Their Coxeter--Burau representations nevertheless exist
independently of that restriction.  The arrangement-group interpretation
of Artin groups is also reviewed in \cite[\S2.1]{MarinKrammer}.
\begin{thebibliography}{99}
\bibitem{Squier}
C.~C.~Squier, \emph{Matrix representations of Artin groups},
Proc. Amer. Math. Soc. \textbf{103} (1988), 49--53.
\url{https://www.maths.ed.ac.uk/~v1ranick/papers/squier2.pdf}.
\bibitem{BQ}
A.~Bapat and H.~Queffelec, \emph{Some remarks about the faithfulness
of the Burau representation of Artin--Tits groups}, arXiv:2409.00144.
\url{https://arxiv.org/abs/2409.00144}.
\bibitem{Damianou}
P.~A.~Damianou, \emph{On the characteristic polynomial of Cartan
matrices and Chebyshev polynomials}, arXiv:1110.6620v4.
\url{https://arxiv.org/abs/1110.6620}.
\bibitem{Marin}
I.~Marin, \emph{Infinitesimal Hecke algebras II}, arXiv:0911.1879,
especially Theorems 1 and 4.
\url{https://arxiv.org/abs/0911.1879}.
\bibitem{FK}
L.~Funar and T.~Kohno, \emph{On Burau representations at roots of
unity}, arXiv:0907.0568.
\url{https://arxiv.org/abs/0907.0568}.
\bibitem{DDM}
J.~Dehaene and B.~De Moor, \emph{Clifford group, stabilizer states,
and linear and quadratic operations over GF(2)},
Phys. Rev. A \textbf{68} (2003), 042318.
\url{https://arxiv.org/abs/quant-ph/0304125}.
\bibitem{Planat}
M.~Planat, \emph{Clifford group dipoles and the enactment of
Weyl/Coxeter group $W(E_8)$ by entangling gates},
General Mathematics Notes \textbf{2}(1) (2011), 96--113.
\url{https://arxiv.org/abs/0904.3691}.
\bibitem{ARW}
M.~Amy, N.~J.~Ross and S.~Wesley,
\emph{A sound and complete equational theory for 3-qubit
Toffoli--Hadamard circuits}, arXiv:2407.11152.
\url{https://arxiv.org/abs/2407.11152}.
\bibitem{Negami}
H.~Negami, \emph{Universal Quantum Gate Sets from Middle Convolution:
Braid Dynamics and Exact Synthesis Costs}, revised working draft,
10 November 2026.
\bibitem{NegamiGates}
H.~Negami, \emph{Quantum gates from the middle convolution of twisted
Burau representations}, arXiv:2610.00293v1 (2026).
\url{https://arxiv.org/abs/2610.00293v1}.
\bibitem{NegamiKZ}
H.~Negami, \emph{Long-Moody construction of braid group representations
and Haraoka's multiplicative middle convolution for KZ-type equations},
arXiv:2503.14840v3 (2026).
\url{https://arxiv.org/abs/2503.14840v3}.
\bibitem{Barenco}
A.~Barenco, C.~H.~Bennett, R.~Cleve, D.~P.~DiVincenzo, N.~Margolus,
P.~Shor, T.~Sleator, J.~A.~Smolin and H.~Weinfurter,
\emph{Elementary gates for quantum computation},
Phys. Rev. A \textbf{52} (1995), 3457--3467.
\url{https://arxiv.org/abs/quant-ph/9503016}.

\bibitem{FLW}
M.~H.~Freedman, M.~J.~Larsen and Z.~Wang,
\emph{The two-eigenvalue problem and density of Jones representation
of braid groups}, Comm. Math. Phys. \textbf{228} (2002), 177--199,
especially Theorems 1.1, 1.4 and 1.6.
\url{https://arxiv.org/abs/math/0103200}.

\bibitem{Etingof}
P.~Etingof, \emph{Lie Groups and Lie Algebras}, MIT OpenCourseWare,
18.745 (Fall 2020), \S\S20--23.
\href{https://ocw.mit.edu/courses/18-745-lie-groups-and-lie-algebras-i-fall-2020/mit18_745_f20_lec_full.pdf}{Lecture notes (PDF)}.

\bibitem{SageArr}
The Sage Developers, \emph{SageMath Reference Manual: Library of
Hyperplane Arrangements}, \texttt{Coxeter}, accessed 6 November 2026.
\href{https://doc.sagemath.org/html/en/reference/discrete_geometry/sage/geometry/hyperplane_arrangement/library.html}{Coxeter arrangement constructor}.

\bibitem{M2B}
Macaulay2, \emph{HyperplaneArrangements package: typeB},
accessed 6 November 2026.
\href{https://macaulay2.com/doc/Macaulay2/share/doc/Macaulay2/HyperplaneArrangements/html/_type__B_lp__Z__Z_cm__Ring_rp.html}{Type B arrangement documentation}.

\bibitem{SageArtin}
\begingroup\raggedright
The Sage Developers, \emph{SageMath Reference Manual: Artin Groups},
\texttt{ArtinGroupElement.burau\_matrix}, accessed 6 November 2026.
\url{https://doc.sagemath.org/html/en/reference/groups/sage/groups/artin.html}.
\par\endgroup

\bibitem{MarinKrammer}
I.~Marin, \emph{Krammer representations for complex braid groups},
arXiv:0711.3096, \S2.1.
\url{https://arxiv.org/abs/0711.3096}.

\bibitem{HR}
T.~Hoge and G.~R\"ohrle, \emph{On supersolvable reflection arrangements},
arXiv:1209.1919v3 (2013), Theorem 1.2.
\url{https://arxiv.org/abs/1209.1919v3}.

\bibitem{AH}
S.~Adachi and K.~Hiroe,
\emph{Middle convolution for affine hyperplane arrangements and Aomoto
realizability}, arXiv:2601.00544v4 (29 September 2026).
\url{https://arxiv.org/abs/2601.00544v4}.

\bibitem{Hiroe}
K.~Hiroe, \emph{Middle convolution for Lie algebra representations},
arXiv:2605.09828v3 (13 October 2026), Theorem 4.7.
\url{https://arxiv.org/abs/2605.09828}.

\bibitem{HN}
K.~Hiroe and H.~Negami, \emph{Long-Moody construction of braid
representations and Katz middle convolution}, arXiv:2303.05770 (2023).
\url{https://arxiv.org/abs/2303.05770}.

\bibitem{CS}
D.~C.~Cohen and A.~I.~Suciu,
\emph{Homology of iterated semidirect products of free groups},
J.~Pure Appl.~Algebra \textbf{126} (1998), 87--120.
\href{https://doi.org/10.1016/S0022-4049(96)00153-3}{doi:10.1016/S0022-4049(96)00153-3}.

\bibitem{Kohno}
T.~Kohno, \emph{Hyperplane arrangements, local system homology and
iterated integrals}, Adv.~Stud.~Pure Math.~\textbf{62} (2012), 157--174.
\url{https://www.ms.u-tokyo.ac.jp/~kohno/papers/kohno_aspm.pdf}.
\bibitem{NegamiSignature}
H.~Negami, \emph{A closed signature formula for the Katz--Long--Moody
Hermitian form}, arXiv:2609.11793 (2026), Theorem 1.1.
\url{https://arxiv.org/abs/2609.11793}.

\bibitem{NegamiFinite}
H.~Negami, \emph{Finite images of braid group representations and
algebraic solutions of KZ-type equations}, arXiv:2609.25805 (2026).
\url{https://arxiv.org/abs/2609.25805}.

\end{thebibliography}
\end{document}
